\documentclass[11pt,a4paper]{article}

\usepackage[utf8]{inputenc}
\usepackage[T1]{fontenc}
\usepackage[spanish,es-noshorthands,es-nodecimaldot,es-tabla]{babel}
\usepackage[a4paper,margin=2.2cm]{geometry}
\usepackage{graphicx}
\usepackage{subcaption}
\usepackage{longtable}
\usepackage{array}
\usepackage{xcolor}
\usepackage{amssymb}
\usepackage{url}
\usepackage[hidelinks]{hyperref}

\renewcommand{\arraystretch}{1.25}
\graphicspath{{figuras/}}

% Texto en rojo = lo que falta escribir. Al terminar no debe quedar ninguno.
\newcommand{\pendiente}[1]{\textcolor{red}{[#1]}}

% Inserta una captura; si la imagen aún no existe, muestra un recuadro para que el documento compile.
\newcommand{\captura}[1]{%
  \IfFileExists{figuras/#1}%
    {\includegraphics[width=\linewidth]{#1}}%
    {\fbox{\parbox[c][4cm][c]{\linewidth}{\centering\pendiente{falta figuras/#1}}}}%
}

\begin{document}

% ------------------------------------------------------------------
% Carátula
% ------------------------------------------------------------------
\begin{titlepage}
	\centering
	\begin{minipage}[c]{0.22\textwidth}\centering
		\includegraphics[height=3.4cm]{logo-uteq.png}
	\end{minipage}\hfill
	\begin{minipage}[c]{0.52\textwidth}\centering
		{\large\bfseries UNIVERSIDAD TÉCNICA ESTATAL DE QUEVEDO}\\[6pt]
		Facultad de Ciencias de la Computación\\
		Carrera de Software
	\end{minipage}\hfill
	\begin{minipage}[c]{0.22\textwidth}\centering
		\includegraphics[height=3.4cm]{logo-fcc.png}
	\end{minipage}

	\vspace*{\fill}
	{\large Asignatura: Ingeniería de Requerimientos (ISR-401)}\\[18pt]
	{\LARGE\bfseries TA-IND-01\\[6pt] IEEE 29148 y tipos de requisitos\\[4pt] de un sistema conocido}\\[10pt]
	{\normalsize Trabajo autónomo individual (actividad sumativa)}
	\vspace*{\fill}

	{\large Estudiante: Vélez Vinces Jostin Gustavo}\\[8pt]
	Cuarto nivel $\cdot$ Paralelo B\\[8pt]
	Docente: Ing. Gleiston Cicerón Guerrero Ulloa, PhD\\[8pt]
	Período académico 2026--2027 SPA
	\vspace*{\fill}

	Quevedo, 11 de noviembre de 2026
\end{titlepage}

% ------------------------------------------------------------------
\section{Sistema elegido}
% ------------------------------------------------------------------
% Un párrafo: propósito, usuarios y contexto de la banca móvil de Banco Pichincha.
La aplicación Banca Móvil del Banco Pichincha es una herramienta digital que permite a sus usuarios realizar operaciones bancarias desde sus teléfonos celulares sin necesidad de acudir a una agencia. A través de ella se pueden consultar saldos y movimientos, realizar transferencias y pagar servicios básicos (Figura~\ref{fig:capturas}). Está dirigida a los clientes del Banco Pichincha que necesitan administrar su dinero mediante herramientas digitales. En Ecuador, se utiliza desde cualquier lugar que cuente con conexión a internet, como el hogar, el trabajo o durante los desplazamientos, permitiendo realizar operaciones bancarias en diferentes momentos del día. De esta manera, facilita el acceso a los servicios bancarios y reduce la necesidad de acudir presencialmente a una agencia.

% Fotografías propias de la pantalla, tomadas con otro teléfono porque la aplicación
% no permite capturas de pantalla. Los datos personales están pixelados.
\begin{figure}[htbp]
	\centering
	\begin{subfigure}[t]{0.185\textwidth}\centering
		\captura{captura-01.png}
		\caption{Inicio de sesión.}\label{fig:login}
	\end{subfigure}\hfill
	\begin{subfigure}[t]{0.185\textwidth}\centering
		\captura{captura-02.png}
		\caption{Menú principal.}\label{fig:menu}
	\end{subfigure}\hfill
	\begin{subfigure}[t]{0.185\textwidth}\centering
		\captura{captura-03.png}
		\caption{Transferencia.}\label{fig:transferencia}
	\end{subfigure}\hfill
	\begin{subfigure}[t]{0.185\textwidth}\centering
		\captura{captura-04.png}
		\caption{Pago de servicios.}\label{fig:servicios}
	\end{subfigure}\hfill
	\begin{subfigure}[t]{0.185\textwidth}\centering
		\captura{captura-05.png}
		\caption{Descuentos y beneficios.}\label{fig:descuentos}
	\end{subfigure}
	\caption{Pantallas de la aplicación de banca móvil de Banco Pichincha (versión 15.0.0). Fotografías propias tomadas con otro teléfono, porque la aplicación no permite capturas de pantalla; los datos personales están pixelados.}
	\label{fig:capturas}
\end{figure}

% ------------------------------------------------------------------
\section{Resumen analítico de ISO/IEC/IEEE 29148:2018}
% ------------------------------------------------------------------
% De 1,5 a 2 páginas, con sus palabras (no traducir). Citar la cláusula exacta de cada idea,
% por ejemplo: «... (cláusula 5.2.5) \cite{iso29148}».

\subsection{Requisito bien formado}
SWEBOK define un requisito de software como una propiedad que algo debe tener para resolver un problema del mundo real \cite{swebok2024}. La norma lo define como un enunciado que expresa una necesidad junto con sus condiciones y restricciones (cláusula 5.2.4). Además, pide que los requisitos estén “bien formados”, porque así es más fácil validarlos con los interesados y comprobar que reflejan lo que de verdad necesitan (cláusula 5.2.4).

Según la norma, un requisito bien formado presenta una o más de estas propiedades: lo debe cumplir el sistema para resolver un problema o lograr un objetivo; tiene condiciones medibles; está limitado por restricciones; describe lo que hace el sistema y no lo que hace el usuario; y se puede verificar (cláusula 5.2.4).

Para escribirlo, la norma pide un sujeto (el sistema), un verbo y la acción o restricción. “Deberá” (\emph{shall}) se usa solo para lo obligatorio, en voz activa y en forma positiva, y se evita “deberá ser capaz de” (cláusula 5.2.4). La norma también diferencia las condiciones, que son atributos medibles del requisito, de las restricciones, que limitan el diseño, como una ley, el presupuesto o una plataforma que ya existe (cláusula 5.2.4).

\textbf{Ejemplo} (pantalla “Ingresa un monto” de la aplicación, Figura~\ref{fig:transferencia}): “La aplicación deberá mostrar el aviso ‘Esta transferencia no tiene costo’ antes de que el cliente confirme una transferencia sin costo”. El sujeto es la aplicación, la acción es obligatoria, tiene una condición y se puede comprobar viendo la pantalla.

\subsection{Características de un requisito individual}
La norma dice que cada requisito, por sí solo, debe cumplir nueve características (cláusula 5.2.5):

\begin{enumerate}
	\item \textbf{Necesario} (\emph{necessary}): si se quita, falta algo que ningún otro requisito cubre, y además sigue vigente (cláusula 5.2.5).
	\item \textbf{Apropiado} (\emph{appropriate}): tiene el nivel de detalle que corresponde y no impone diseño sin necesidad (cláusula 5.2.5). Ejemplo: “La aplicación deberá permitir el ingreso mediante huella dactilar” es apropiado. “La app deberá usar la librería biométrica X” no lo es, porque ya decide el diseño.
	\item \textbf{Inequívoco} (\emph{unambiguous}): se entiende de una sola forma (cláusula 5.2.5). Si un requisito dice “mostrar el saldo rápidamente”, el programador y el banco pueden entender tiempos distintos.
	\item \textbf{Completo} (\emph{complete}): se entiende sin buscar información en otro lado (cláusula 5.2.5). “La app deberá permitir transferir” queda incompleto si no dice desde qué cuenta, aunque la pantalla de la aplicación sí tiene el campo “Desde” (Figura~\ref{fig:transferencia}).
	\item \textbf{Singular} (\emph{singular}): pide una sola capacidad o restricción (cláusula 5.2.5). Ejemplo: “La app permitirá ver el saldo, ocultarlo con el ícono del ojo y transferir dinero” se debe separar en tres requisitos de la pantalla “Mis productos” (Figura~\ref{fig:menu}).
	\item \textbf{Factible} (\emph{feasible}): se puede hacer con el costo, el tiempo y la tecnología disponibles (cláusula 5.2.5). Pedir que una transferencia llegue en 0 segundos no es factible.
	\item \textbf{Verificable} (\emph{verifiable}): se puede demostrar que se cumplió, y es más fácil cuando es medible (cláusula 5.2.5). SWEBOK también considera la verificabilidad una propiedad esencial de los requisitos \cite{swebok2024}. Si un requisito no se puede probar, no sirve, porque nadie puede saber si se cumplió. Ejemplo: “La app deberá mostrar promociones relevantes” no es verificable. En cambio, “La tarjeta de una promoción en la pantalla ‘Descuentos y beneficios’ deberá mostrar: establecimiento, descuento, cobertura, fecha de vigencia” se puede revisar dato por dato (Figura~\ref{fig:descuentos}).
	\item \textbf{Correcto} (\emph{correct}): representa bien la necesidad de la que salió (cláusula 5.2.5). Si el banco quiere avisar que la transferencia es gratis y el requisito pide “mostrar el costo de la transferencia”, el requisito no es correcto.
	\item \textbf{Conforme} (\emph{conforming}): sigue la plantilla y el estilo de redacción aprobados (cláusula 5.2.5). Por ejemplo, todos los requisitos de este trabajo usan “deberá” como verbo obligatorio, con un sujeto, una acción y una condición. Si un requisito estuviera escrito con otra estructura, por ejemplo con “debe” o en voz pasiva, no sería conforme.
\end{enumerate}

\subsection{Características de un conjunto de requisitos}
La norma aclara que algunas cualidades se revisan en el conjunto y no en cada requisito por separado (cláusula 5.2.6):

\begin{itemize}
	\item \textbf{Completo} (\emph{complete}): el conjunto describe todo lo necesario y no deja pendientes “por definir” (TBD), “por especificar” (TBS) o “por resolver” (TBR) (cláusula 5.2.6).
	\item \textbf{Consistente} (\emph{consistent}): los requisitos no se repiten ni se contradicen, usan las mismas unidades y cada término significa siempre lo mismo (cláusula 5.2.6). Ejemplo: en la aplicación, la función QR aparece como “Pagar con QR”, “Ver mi QR” y “Escanear o mostrar QR” (Figuras~\ref{fig:login} y~\ref{fig:menu}). La especificación tiene que definir cada término para no llamar de formas distintas a una misma función.
	\item \textbf{Factible} (\emph{feasible}): se puede hacer el conjunto entero, y es asequible (cláusula 5.2.6).
	\item \textbf{Comprensible} (\emph{comprehensible}): se entiende qué se espera y cómo se relaciona con el sistema mayor (cláusula 5.2.6).
	\item \textbf{Validable} (\emph{able to be validated}): si se cumple el conjunto, se satisfacen de verdad las necesidades dentro de los límites técnicos, de costo y legales (cláusula 5.2.6).
\end{itemize}

La norma advierte que revisar esto a tiempo evita que los requisitos crezcan sin control (\emph{requirements creep}) (cláusula 5.2.6).

\subsection{Términos ambiguos que la norma prohíbe}
La norma pide que los requisitos digan “qué” se necesita y no “cómo”. También pide evitar los términos vagos, porque no dejan verificar o se pueden entender de varias formas (cláusula 5.2.7). Los tipos son:

\begin{itemize}
	\item superlativos (“el mejor”);
	\item lenguaje subjetivo (“fácil de usar”);
	\item pronombres vagos (“esto”, “eso”);
	\item adverbios y adjetivos ambiguos (“casi siempre”, “mínimo”) y conectores como “y/o”;
	\item términos abiertos (“pero no limitado a”);
	\item comparativos (“mejor que”);
	\item escapatorias (“si es posible”, “según corresponda”);
	\item términos de totalidad (“todos”, “siempre”, “nunca”);
	\item referencias incompletas, sin versión ni fecha (cláusula 5.2.7).
\end{itemize}

\textbf{Ejemplo} (pantalla “Pago de servicios”, Figura~\ref{fig:servicios}): “La app deberá permitir pagar todos los servicios de forma fácil y, si es posible, programarlos” tiene un término de totalidad, lenguaje subjetivo y una escapatoria. Se corrige con dos requisitos: “La aplicación deberá mostrar las siguientes categorías: Agua Potable, Electricidad, SRI, IESS, Venta por Catálogo, Almacenes” y “La aplicación deberá permitir crear una programación de pago desde la opción ‘Crear programación’”. Por último, la norma pide documentar los supuestos en los atributos del requisito (cláusula 5.2.7).

% ------------------------------------------------------------------
\section{Glosario}
% ------------------------------------------------------------------
% 10 términos de la cláusula 3.1, definición parafraseada y un ejemplo propio en la app.
{\small
\begin{longtable}{|>{\raggedright\arraybackslash}p{0.5cm}|>{\raggedright\arraybackslash}p{2.6cm}|>{\raggedright\arraybackslash}p{1.4cm}|>{\raggedright\arraybackslash}p{5.4cm}|>{\raggedright\arraybackslash}p{5.4cm}|}
\caption{Glosario de términos de la cláusula 3.1 de ISO/IEC/IEEE 29148:2018 \cite{iso29148}.}\label{tab:glosario}\\
\hline
\textbf{N.º} & \textbf{Término} & \textbf{Cláusula} & \textbf{Definición parafraseada} & \textbf{Ejemplo en la banca móvil} \\ \hline
\endfirsthead
\hline
\textbf{N.º} & \textbf{Término} & \textbf{Cláusula} & \textbf{Definición parafraseada} & \textbf{Ejemplo en la banca móvil} \\ \hline
\endhead
1 & Requisito (\emph{requirement}) & 3.1.19 & Enunciado que expresa una necesidad junto con sus condiciones y restricciones. Siempre se refiere a un sistema, software o servicio. & “La aplicación deberá mostrar el aviso ‘Esta transferencia no tiene costo’ antes de que el cliente confirme una transferencia sin costo.” \\ \hline
2 & Interesado (\emph{stakeholder}) & 3.1.28 & Persona u organización que tiene un derecho, un interés o un reclamo sobre el sistema. Incluye usuarios, desarrolladores, clientes, operadores y entes reguladores. & El cliente, Banco Pichincha, el equipo que desarrolla la aplicación, la Superintendencia de Bancos y los comercios que aparecen en “Descuentos y beneficios”. \\ \hline
3 & Usuario (\emph{user}) & 3.1.35 & Persona o grupo que interactúa con el sistema o se beneficia de él mientras se usa. Es un tipo de interesado. & El cliente que entra con su huella y paga la luz en “Pago de servicios”. También quien recibe una transferencia se beneficia del sistema aunque no lo use en ese momento. \\ \hline
4 & Condición (\emph{condition}) & 3.1.6 & Atributo medible que indica en qué circunstancia o evento aplica un requisito. & En el requisito del aviso, la condición es “antes de que el cliente confirme la transferencia”. \\ \hline
5 & Restricción (\emph{constraint}) & 3.1.7 & Limitación impuesta desde fuera sobre el sistema, su diseño o el proceso con que se desarrolla. No la decide el equipo, se le impone. & La LOPDP obliga a la aplicación a tratar los datos personales del cliente según la ley, aunque al equipo le resultara más fácil hacerlo de otra forma. \\ \hline
6 & Atributo (\emph{attribute}) & 3.1.2 & Propiedad de algo que se puede distinguir por cantidad o por cualidad. & Los atributos de un requisito, como su ID (RQ-01), su tipo y su prioridad. En la aplicación, el “Saldo disponible” es un atributo de la cuenta. \\ \hline
7 & Verificación (\emph{verification}) & 3.1.37 & Confirmar con evidencia que se cumplió lo que estaba especificado. Responde a “¿lo construimos bien?”. & Un \emph{tester} abre “Ingresa un monto” y comprueba que el aviso aparece con el texto exacto que pide el requisito. \\ \hline
8 & Validación (\emph{validation}) & 3.1.36 & Confirmar con evidencia que el sistema sirve para el uso previsto. Responde a “¿construimos el sistema correcto?”. & Se pide a clientes reales que paguen su planilla de agua desde la aplicación. Si no logran hacerlo sin ayuda, la aplicación no queda validada, aunque cumpla cada requisito escrito. \\ \hline
9 & Trazabilidad de requisitos (\emph{requirements traceability}) & 3.1.23 & Registrar de dónde viene un requisito (hacia arriba, el requisito “padre”) y en qué otros se divide (hacia abajo, los requisitos “hijos”). & La necesidad “pagar servicios desde el celular” (padre) se divide en RQ-05, que muestra las categorías, y RQ-06, que permite crear una programación de pago (hijos). \\ \hline
10 & Sistema de interés (\emph{system-of-interest}) & 3.1.32 & El sistema cuyo ciclo de vida se está analizando en el trabajo. & La aplicación móvil de Banco Pichincha (versión 15.0.0). La Banca Web, Deuna y el sistema del SRI quedan fuera, aunque la aplicación se conecte con ellos. \\ \hline
\end{longtable}
}

% ------------------------------------------------------------------
\section{Requisitos del sistema}
% ------------------------------------------------------------------
% 15 requisitos, al menos 5 tipos distintos (tipos de la cláusula 5.2.8.3) y al menos 5 de calidad
% con característica y subcaracterística de ISO/IEC 25010:2023 y una métrica.
% Para los que no son de calidad, en «Característica» y «Métrica» poner «N/A».
Se formularon 15 requisitos para la aplicación móvil de Banco Pichincha, basados principalmente en las pantallas de la Figura~\ref{fig:capturas}. El tipo de cada requisito se asignó según la cláusula 5.2.8.3 de la norma ISO/IEC/IEEE 29148:2018 \cite{iso29148}, y los requisitos de calidad se clasificaron según las características y subcaracterísticas de ISO/IEC 25010:2023 \cite{iso25010}. Cada requisito incluye una métrica, cuando corresponde, y un criterio de aceptación que permite verificar su cumplimiento (Tabla~\ref{tab:requisitos}).

{\footnotesize
\setlength{\tabcolsep}{3pt}
\begin{longtable}{|>{\raggedright\arraybackslash}p{0.95cm}|>{\raggedright\arraybackslash}p{4.7cm}|>{\raggedright\arraybackslash}p{1.6cm}|>{\raggedright\arraybackslash}p{2.4cm}|>{\raggedright\arraybackslash}p{2.1cm}|>{\raggedright\arraybackslash}p{3.7cm}|}
\caption{Requisitos de la banca móvil de Banco Pichincha clasificados según \cite{iso29148} e \cite{iso25010}.}\label{tab:requisitos}\\
\hline
\textbf{ID} & \textbf{Enunciado} & \textbf{Tipo} & \textbf{Característica ISO/IEC 25010} & \textbf{Métrica} & \textbf{Criterio de aceptación} \\ \hline
\endfirsthead
\hline
\textbf{ID} & \textbf{Enunciado} & \textbf{Tipo} & \textbf{Característica ISO/IEC 25010} & \textbf{Métrica} & \textbf{Criterio de aceptación} \\ \hline
\endhead
RQ-01 & La aplicación deberá permitir el ingreso del cliente con la huella dactilar registrada en su teléfono. & Funcional & N/A & N/A & Con una huella registrada, el cliente pulsa “Huella / Face ID” y entra a “Mis productos”. \\ \hline
RQ-02 & La aplicación deberá impedir la captura de pantalla en la pantalla de inicio de sesión. & Calidad & Seguridad – Confidencialidad & Capturas legibles guardadas & Se intentan 10 capturas en el inicio de sesión y la galería guarda 0 imágenes legibles. \\ \hline
RQ-03 & La aplicación deberá ocultar el saldo disponible cuando el cliente pulse el ícono del ojo en “Mis productos”. & Funcional & N/A & N/A & Al pulsar el ícono, el saldo se muestra como “**.**”. \\ \hline
RQ-04 & La aplicación deberá mostrar el aviso “Esta transferencia no tiene costo” en “Ingresa un monto” antes de que el cliente confirme una transferencia sin costo. & Funcional & N/A & N/A & En una transferencia sin costo, el aviso aparece antes del botón de confirmación. \\ \hline
RQ-05 & La aplicación deberá mostrar en “Pago de servicios” las siguientes categorías: Agua Potable, Electricidad, SRI, IESS, Venta por Catálogo, Almacenes. & Funcional & N/A & N/A & Al abrir “Pago de servicios” se ven las seis categorías. \\ \hline
RQ-06 & La aplicación deberá permitir crear una programación de pago desde la opción “Crear programación”. & Funcional & N/A & N/A & El cliente crea una programación con fecha y servicio, y esta queda guardada en la pantalla. \\ \hline
RQ-07 & La aplicación deberá confirmar una transferencia entre cuentas de Banco Pichincha en un máximo de 5~s en el 95~\% de los casos, con conexión 4G. & Calidad & Eficiencia de desempeño – Comportamiento temporal & Tiempo de respuesta (s) & De 100 transferencias con 4G, al menos 95 se confirman en $\leq$~5~s. \\ \hline
RQ-08 & La aplicación deberá mostrar, antes de ejecutar una transferencia, una pantalla de confirmación con el nombre del beneficiario, el monto y la cuenta de origen. & Factores humanos & Capacidad de interacción – Protección contra errores del usuario & Transferencias ejecutadas sin confirmación (\%) & De 20 transferencias de prueba, el 0~\% se ejecuta sin la pantalla de confirmación. \\ \hline
RQ-09 & La aplicación deberá consultar al SRI el valor de la matrícula vehicular según la placa ingresada en el campo “PLACA RAMV CPN”, cuando el cliente pulse “Validar”. & Interfaz & Compatibilidad – Interoperabilidad & Valores que coinciden con el SRI (\%) & Con 10 placas de prueba, el valor coincide con el del SRI en 10 de 10 consultas. \\ \hline
RQ-10 & La aplicación deberá filtrar las promociones de “Descuentos y beneficios” según la ciudad que elija el cliente. & Funcional & N/A & N/A & Al elegir “Quevedo”, solo aparecen promociones válidas en Quevedo. \\ \hline
RQ-11 & La aplicación deberá estar disponible el 99,5~\% del tiempo mensual, sin contar mantenimientos avisados con 48~h de anticipación. & Calidad & Fiabilidad – Disponibilidad & Disponibilidad mensual (\%) & Un monitoreo de 30 días registra $\geq$~99,5~\% de tiempo en línea. \\ \hline
RQ-12 & La aplicación deberá permitir que un cliente que nunca la ha usado pague una planilla de electricidad en un máximo de 3~min sin ayuda. & Usabilidad & Capacidad de interacción – Aprendizaje & Tiempo de tarea (min) y tasa de éxito (\%) & De 10 clientes nuevos, al menos 9 pagan en $\leq$~3~min sin ayuda. \\ \hline
RQ-13 & La aplicación deberá bloquear el ingreso con PIN durante 30~min tras 3 intentos fallidos consecutivos. & Calidad & Seguridad – Autenticidad & Intentos antes del bloqueo; duración (min) & Tras 3 PIN incorrectos, el cuarto intento se rechaza y el PIN vuelve a funcionar a los 30~min. \\ \hline
RQ-14 & El tratamiento de los datos personales del cliente en la aplicación deberá cumplir la Ley Orgánica de Protección de Datos Personales (Registro Oficial, Quinto Suplemento N.º 459, 26 de mayo de 2021). & Restricción & N/A & N/A & Una auditoría de cumplimiento de la LOPDP no reporta incumplimientos. \\ \hline
RQ-15 & El banco deberá publicar las notas de versión de una actualización de la aplicación antes de su liberación en la tienda de aplicaciones. & Proceso & N/A & N/A & El día de la liberación, la ficha de la tienda muestra las notas de esa versión. \\ \hline
\end{longtable}
}

% ------------------------------------------------------------------
\section{Autoevaluación}
% ------------------------------------------------------------------
% Revisar cada requisito contra las nueve características de 5.2.5 y los términos de 5.2.7.
% En la tabla van solo los que fallaron: la versión anterior, qué falló y la corregida.
% Si se corrige uno, en la Tabla~\ref{tab:requisitos} va la versión corregida.
\pendiente{Párrafo corto: cómo se revisaron los requisitos y cuántos se corrigieron.}

{\small
\begin{longtable}{|>{\raggedright\arraybackslash}p{0.9cm}|>{\raggedright\arraybackslash}p{4.6cm}|>{\raggedright\arraybackslash}p{3.4cm}|>{\raggedright\arraybackslash}p{5.6cm}|}
\caption{Requisitos corregidos en la autoevaluación.}\label{tab:autoevaluacion}\\
\hline
\textbf{ID} & \textbf{Versión anterior} & \textbf{Qué falló (5.2.5 / 5.2.7)} & \textbf{Versión corregida} \\ \hline
\endfirsthead
\hline
\textbf{ID} & \textbf{Versión anterior} & \textbf{Qué falló (5.2.5 / 5.2.7)} & \textbf{Versión corregida} \\ \hline
\endhead
\pendiente{RQ-xx} & \pendiente{versión anterior} & \pendiente{característica o término} & \pendiente{versión corregida} \\ \hline
\end{longtable}
}

% ------------------------------------------------------------------
\section{Reflexión}
% ------------------------------------------------------------------
% Media página: qué tipo de requisito costó más formular y por qué.
\pendiente{Escribir la reflexión.}

% ------------------------------------------------------------------
\section{Declaración de uso de inteligencia artificial}
% ------------------------------------------------------------------
% Completar al final con lo que realmente se hizo con la IA.
\begin{description}
	\item[Herramienta:] Claude, asistente de inteligencia artificial de Anthropic.
	\item[Versión:] Claude Opus 5.5.
	\item[Propósito:] preparar la plantilla \LaTeX{} del documento, el archivo README.md y el archivo de referencias; ubicar en la norma las páginas de las cláusulas 1, 3 y 5.2; \pendiente{completar con el resto del uso real}.
	\item[Cómo se verificó lo obtenido:] \pendiente{completar}.
\end{description}

\bibliographystyle{IEEEtran}
\bibliography{referencias}

\end{document}
